\subsection{Hensel's Lemma}

In a topological ring A, an element x is called topologically nilpotent if 0 is a limit of the sequence \((x^{n})_{n\geqslant0}\) . If A is a linearly topologized commutative ring, to say that x is topologically nilpotent means that for every open ideal \(\mathfrak{S}\) of A the canonical image of x in A/\(\mathfrak{S}\) is a nilpotent element of that ring. If \(r_{\mathfrak{S}}\) is the nil-radical of A/\(\mathfrak{S}\), clearly \((r_{\mathfrak{S}})\) is an inverse system of subsets and the set t of topological nilpotent elements of A is the inverse image of \(r=\lim_{\leftarrow \mathfrak{S}}r_{\mathfrak{S}}\) under the

canonical homomorphism \(\mathbf{A} \to \lim_{\leftarrow \mathfrak{S}} \mathbf{A} / \mathfrak{S}\) ; it is therefore a closed ideal of \(\mathbf{A}\) . If also \(\mathbf{A}\) is Hausdorff and complete, this ideal is contained in the Jacobson radical of \(\mathbf{A}\) and, for an element \(x \in \mathbf{A}\) to be invertible, it is necessary and sufficient that its class mod. \(t\) be invertible in \(\mathbf{A} / t\) (§ 2, no. 13, Lemma 3).

Note that if A is a ring and m a two-sided ideal of A, the elements of m are topologically nilpotent with respect to the m-adic topology.

THEOREM 1 (Hensel). Let A be a complete Hausdorff linearly topologized commutative ring. Let m be a closed ideal of A whose elements are topologically nilpotent. Let B = A/m be the quotient topological ring and \(\varphi: A \to B\) the canonical mapping. Let R be a restricted formal power series in A\{X\}, \(\overline{\mathrm{P}}\) a monic polynomial in B{[}X{]} and \(\overline{\mathrm{Q}}\) a restricted formal power series in B\{X\}. Suppose that \(\overline{\varphi}(\mathrm{R}) = \overline{\mathrm{P}}\) . \(\overline{\mathrm{Q}}\) and that \(\overline{\mathrm{P}}\) and \(\overline{\mathrm{Q}}\) are strongly relatively prime in B\{X\}. Then there exists a unique ordered pair (P, Q) consisting of a monic polynomial \(\mathbf{P} \in \mathbf{A}[\mathbf{X}]\) and a restricted formal power series \(\mathbf{Q} \in \mathbf{A}\{\mathbf{X}\}\) such that

\[
\mathrm{R} = \mathrm{P}. \mathrm{Q}, \quad \overline {{{\varphi}}} (\mathrm{P}) = \overline {{{\mathrm{P}}}}, \quad \overline {{{\varphi}}} (\mathrm{Q}) = \overline {{{\mathrm{Q}}}}.\tag{4}
\]

Moreover, P and Q are strongly relatively prime in A\{X\} and, if R is a polynomial, so is Q.

The proof is divided into four steps. In the first three we assume that A is discrete, in which case R and \(\overline{Q}\) are polynomials.

\textbf{(1) \(\mathfrak{m}^2 = 0\)}\par

Let S, T be two polynomials of A{[}X{]} such that S is monic and \(\overline{\varphi}(S) = \overline{P}\) , \(\overline{\varphi}(T) = \overline{Q}\) ; Proposition 2 of no. 1 shows that S and T are strongly relatively prime; hence (no. 1, Proposition 1) there exists a unique ordered pair of polynomials (S', T') of A{[}X{]} such that

\[
\mathrm{R} - \mathrm{ST} = \mathrm{ST} ^ {\prime} + \mathrm{TS} ^ {\prime} \quad \text { and } \quad \deg (\mathrm{S} ^ {\prime}) <   \deg (\mathrm{S}) = \deg (\overline {{\mathrm{P}}}).\tag{5}
\]

The polynomials \(P = S + TS'\) , \(Q = T + T'\) are then solutions to the problem; then in fact

\[
\overline {{{\mathbf {P}}}} \cdot \overline {{{\varphi}}} (\mathbf {T} ^ {\prime}) + \overline {{{\mathbf {Q}}}} \cdot \overline {{{\varphi}}} (\mathbf {S} ^ {\prime}) = \overline {{{\varphi}}} (\mathbf {S T} ^ {\prime} + \mathbf {T S} ^ {\prime}) = \overline {{{\varphi}}} (\mathbf {R} - \mathbf {S T}) = 0.\tag{6}
\]

As \(\overline{\mathbf{P}}\) is monic, \(\overline{\mathbf{P}}\) and \(\overline{\mathbf{Q}}\) strongly relatively prime and

\[
\deg (\overline {{{\varphi}}} (S ^ {\prime})) <   \deg (\overline {{{\mathbf {P}}}}),
\]

Proposition 1 of no. 1 shows that \(\bar{\varphi}(\mathbf{S}') = \bar{\varphi}(\mathbf{T}') = 0\) , in other words the coefficients of \(\mathbf{S}'\) and \(\mathbf{T}'\) belong to \(\mathfrak{m}\) and the relation \(\mathfrak{m}^2 = 0\) gives

\[
\mathrm{PQ} = \mathrm{ST} + \mathrm{ST} ^ {\prime} + \mathrm{TS} ^ {\prime} = \mathrm{R},
\]

which satisfies relation (4). Since \(\overline{\varphi}(\mathbf{P}) = \overline{\mathbf{P}}\) and \(\overline{\varphi}(\mathbf{Q}) = \overline{\mathbf{Q}}\) , \(\mathbf{P}\) and \(\mathbf{Q}\) are strongly relatively prime (no. 1, Proposition 2); finally, if \(\mathbf{P}_1\) and \(\mathbf{Q}_1\) are two other polynomials in A{[}X{]} satisfying (4) and such that \(\mathbf{P}_1\) is monic, then necessarily, setting \(S_1' = P_1 - S\) , \(T_1' = Q_1 - T\) , \(\deg(S_1') < \deg(S)\) and \(R - ST = ST_1' + TS_1'\) since \(S_1'\) and \(T_1'\) have their coefficients in \(m\) ; but Proposition 1 then proves that \(S' = S_1'\) and \(T' = T_1'\) , which proves the uniqueness of the ordered pair \((P, Q)\) .

\textbf{(2) \(\mathfrak{m}\) is nilpotent}\par

Let \(n\) be the smallest integer such that \(m^n = 0\) and let us argue by induction on \(n > 2\) , the theorem having been shown for \(n = 2\) . Let \(A' = A / m^{n-1}\) , \(m' = m / m^{n-1}\) ; as \(m'^{n-1} = 0\) , there exists a unique ordered pair \((P', Q')\) of polynomials in \(A'[X]\) such that \(P'\) is monic, \(R' = P'Q'\) , \(\bar{\psi}(P') = \overline{P}\) and \(\bar{\psi}(Q') = \overline{Q}\) , where \(\psi\) denotes the canonical homomorphism \(A' \to A'/m' = B\) , \(\theta\) the canonical homomorphism \(A \to A'\) and \(R' = \bar{\theta}(R)\) . On the other hand, as \((\mathfrak{m}^{n-1})^{2}=0\) , there exists a unique ordered pair \((\mathrm{P},\mathrm{Q})\) of polynomials in A{[}X{]} such that P is monic and \(R=PQ,\bar{\theta}(P)=\overline{P},\bar{\theta}(Q)=\overline{Q}\) ; as \(\phi=\psi\circ\theta\) , this shows the existence and uniqueness of P and Q satisfying (4); moreover \(P'\) and \(Q'\) are strongly relatively prime by the induction hypothesis and hence so are P and Q.

\textbf{(3) A is discrete}\par

Note that in this case m is no longer necessarily nilpotent, but it is always a nilideal by hypothesis. Let \(P_{0}\) , \(Q_{0}\) be two polynomials of A{[}X{]} such that \(\overline{\varphi}(P_{0}) = \overline{P}\) , \(\overline{\varphi}(Q_{0}) = \overline{Q}\) and \(P_{0}\) is monic. Let us consider the ideal n of A generated by the coefficients of \(R - P_{0}Q_{0}\) ; it is finitely generated and contained in m, hence it is nilpotent (Chapter II, §2, no. 6, Proposition 15) and by definition, if \(\psi: A \to A/n\) is the canonical mapping, then \(\overline{\psi}(R) = \overline{\psi}(P_{0})\overline{\psi}(Q_{0})\) . Moreover, \(\overline{\psi}(P_{0})\) and \(\overline{\psi}(Q_{0})\) are strongly relatively prime, as follows from the hypothesis on P and Q and Proposition 2 of no. 1 applied to the canonical homomorphism \(A/n \to A/m\) . By virtue of case (2), there therefore exists an ordered pair \((P, Q)\) of polynomials in A{[}X{]} such that P is monic and relations (4) hold. The fact that P and Q are strongly relatively prime implies also here that P and Q are strongly relatively prime in A{[}X{]} by virtue of no. 1, Proposition 2, for m is contained in the Jacobson radical of A. Suppose finally that \(P_{1}\) , \(Q_{1}\) are two polynomials in A{[}X{]} satisfying (4) and such that \(P_{1}\) is monic and let \(n_{1}\) be the finitely generated ideal of A generated by the coefficients of \(P - P_{1}\) and the coefficients of \(Q - Q_{1}\) ; as \(n_{1}\) is contained in m, it is nilpotent and, if \(\psi_{1}: A \to A/n_{1}\) is the canonical mapping, \(\overline{\psi}_{1}(P) = \overline{\psi}_{1}(P_{1})\) and

\[
\overline {{{{\psi}}}} _ {1} (\mathbf {Q}) = \overline {{{{\psi}}}} _ {1} (\mathbf {Q} _ {1});
\]

the uniqueness property for case (2) therefore implies \(\mathrm{P} = \mathrm{P}_1, \mathrm{Q} = \mathrm{Q}_1\) .

\textbf{(4) General case}\par

Let \(\mathcal{B}\) be a fundamental system of neighbourhoods of 0 in A consisting of ideals of A. For all \(\mathfrak{S} \in \mathcal{B}\) , let \(f_{\mathfrak{S}}\) be the canonical mapping \(\mathrm{A} \to \mathrm{A} / \mathfrak{S}\) , \(\phi_{\mathfrak{S}}\) the canonical mapping

\[
\mathrm{A} / \mathfrak{S} \mapsto (\mathrm{A} / \mathfrak{S}) / ((\mathfrak {m} + \mathfrak{S}) / \mathfrak{S}) = \mathrm{A} / (\mathfrak {m} + \mathfrak{S}),
\]

\(g_{\mathfrak{S}}\) the canonical mapping \(\mathrm{B} = \mathrm{A} / \mathfrak{m}\to \mathrm{A} /(\mathfrak{m} + \mathfrak{S})\) and write \(\mathrm{R}_{\mathfrak{S}} = \bar{f}_{\mathfrak{S}}(\mathrm{R})\) , \(\overline{\mathrm{P}}_{\mathfrak{S}} = \bar{g}_{\mathfrak{S}}(\overline{\mathrm{P}}),\overline{\mathrm{Q}}_{\mathfrak{S}} = \bar{g}_{\mathfrak{S}}(\overline{\mathrm{Q}})\) . As each ring \(\mathrm{A} / \mathfrak{S}\) is discrete, case (3) can be applied to it and we see that there exists a unique ordered pair \((\mathrm{P}_{\mathfrak{S}},\mathrm{Q}_{\mathfrak{S}})\) of polynomials in \((\mathrm{A} / \mathfrak{S})[\mathrm{X}]\) such that \(\mathrm{P}_{\mathfrak{S}}\) is monic and \(\mathrm{R}_{\mathfrak{S}} = \mathrm{P}_{\mathfrak{S}}\mathrm{Q}_{\mathfrak{S}},\overline{\varphi}_{\mathfrak{S}}(\mathrm{P}_{\mathfrak{S}}) = \overline{\mathrm{P}}_{\mathfrak{S}},\overline{\varphi}_{\mathfrak{S}}(\mathrm{Q}_{\mathfrak{S}}) = \overline{\mathrm{Q}}_{\mathfrak{S}}\) . The uniqueness of this ordered pair implies that, if \(\mathfrak{S}'\subset \mathfrak{S},\mathfrak{S}'\in \mathcal{B}\) and \(f_{\mathfrak{S}\mathfrak{S}'}:\mathrm{A} / \mathfrak{S}'\rightarrow \mathrm{A} / \mathfrak{S}\) is the canonical mapping, then \(\mathrm{P}_{\mathfrak{S}} = \bar{f}_{\mathfrak{S}\mathfrak{S}'}(\mathrm{P}_{\mathfrak{S}'})\) , \(\mathrm{Q}_{\mathfrak{S}} = \bar{f}_{\mathfrak{S}\mathfrak{S}'}(\mathrm{Q}_{\mathfrak{S}'})\) . Then it follows from the canonical identification of \(\mathrm{A}\{\mathrm{X}\}\) with \(\lim (\mathrm{A} /\mathfrak{S})[\mathrm{X}]\)

\[
\varprojlim_ {\mathfrak {S}} (\mathrm{A} / \mathfrak {S}) [ X ]
\]

(no. 2, Proposition 3) that there exists \(\mathbf{P} \in \mathbf{A}\{\mathbf{X}\}\) and \(\mathbf{Q} \in \mathbf{A}\{\mathbf{X}\}\) such that \(\mathbf{R} = \mathbf{PQ}\) and \(\bar{f}_{\mathfrak{S}}(\mathbf{P}) = \mathbf{P}_{\mathfrak{S}}, \bar{f}_{\mathfrak{S}}(\mathbf{Q}) = \mathbf{Q}_{\mathfrak{S}}\) for all \(\mathfrak{S} \in \mathcal{B}\) . Moreover,

\[
\bar {g} _ {\mathfrak {S}} (\overline {{{\mathbf {P}}}} - \overline {{{\varphi}}} (\mathbf {P})) = 0, \quad \bar {g} _ {\mathfrak {S}} (\overline {{{\mathbf {Q}}}} - \overline {{{\varphi}}} (\mathbf {Q})) = 0
\]

for all \(\mathfrak{S} \in \mathcal{B}\) , which means that for all \(\mathfrak{S} \in \mathcal{B}\) the coefficients of \(\overline{\mathbf{P}} - \overline{\varphi}(\mathbf{P})\) and \(\overline{\mathbf{Q}} - \overline{\varphi}(\mathbf{Q})\) all belong to \((\mathfrak{m} + \mathfrak{S})/\mathfrak{m}\) . But, as \(\mathfrak{m}\) is closed in \(\mathbf{A}, \bigcap_{\mathfrak{S}} (\mathfrak{m} + \mathfrak{S}) = \mathfrak{m}\) , whence \(\overline{\mathbf{P}} = \overline{\varphi}(\mathbf{P}), \overline{\mathbf{Q}} = \overline{\varphi}(\mathbf{Q})\) and \(\mathbf{P}\) and \(\mathbf{Q}\) then certainly satisfy (4); moreover, as the \(\mathbf{P}_{\mathfrak{S}}\) are monic and of the same degree, the restricted formal power series \(\mathbf{P}\) is a monic polynomial. If \((\mathbf{P}', \mathbf{Q}')\) were another ordered pair satisfying (4) and such that \(\mathbf{P}'\) is a monic polynomial, we would deduce that \(\mathrm{R}_{\mathfrak{S}} = f_{\mathfrak{S}}(\mathrm{P}') f_{\mathfrak{S}}(\mathrm{Q}')\) , \(\overline{\varphi}_{\mathfrak{S}}(f_{\mathfrak{S}}(\mathrm{P}')) = \overline{\mathrm{P}}_{\mathfrak{S}}\) and \(\overline{\varphi}_{\mathfrak{S}}(f_{\mathfrak{S}}(\mathrm{Q}')) = \overline{\mathrm{Q}}_{\mathfrak{S}}\) and by the uniqueness in case (3) \(f_{\mathfrak{S}}(\mathrm{P}') = \mathrm{P}_{\mathfrak{S}}\) , \(f_{\mathfrak{S}}(\mathrm{Q}') = \mathrm{Q}_{\mathfrak{S}}\) for all \(\mathfrak{S} \in \mathcal{B}\) , which implies that \(\mathrm{P} = \mathrm{P}'\) and \(\mathrm{Q} = \mathrm{Q}'\) . Let us show finally that \(\mathbf{P}\) and \(\mathbf{Q}\) are strongly relatively prime; by virtue of case (3) and Proposition 1 of no. 1, for all \(\mathfrak{S} \in \mathcal{B}\) , there exists a unique ordered pair \((\mathrm{S}_{\mathfrak{S}}, \mathrm{T}_{\mathfrak{S}})\) of polynomials in \((\mathrm{A}/\mathfrak{S})[\mathrm{X}]\) such that

\[
1 = \mathrm{P} _ {\mathfrak {S}} \mathrm{S} _ {\mathfrak {S}} + \mathrm{Q} _ {\mathfrak {S}} \mathrm{T} _ {\mathfrak {S}} \quad \text { and } \quad \deg (\mathrm{T} _ {\mathfrak {S}}) <   \deg (\mathrm{P} _ {\mathfrak {S}}) = \deg (\bar {\mathrm{P}}).\tag{7}
\]

The uniqueness of this ordered pair shows immediately that, for \(\mathfrak{S}' \in \mathcal{B}\) , \(\mathfrak{S}' \subset \mathfrak{S}\) , \(S_{\mathfrak{S}} = \bar{f}_{\mathfrak{S}\mathfrak{S}'}(S_{\mathfrak{S}'})\) , \(T_{\mathfrak{S}} = \bar{f}_{\mathfrak{S}\mathfrak{S}'}(T_{\mathfrak{S}'})\) ; taking account of no. 2, Proposition 3, we conclude that there exist two restricted formal power series \(S\) , \(T\) of \(A\{X\}\) such that \(S_{\mathfrak{S}} = \bar{f}_{\mathfrak{S}}(S)\) , \(T_{\mathfrak{S}} = \bar{f}_{\mathfrak{S}}(T)\) and \(1 = PS + QT\) .

It remains to verify that, if R is a polynomial, so is Q. Now, the \(Q_{\mathfrak{S}}\) are polynomials by construction and, as \(P_{\mathfrak{S}}\) is monic, the relation \(R_{\mathfrak{S}} = P_{\mathfrak{S}}Q_{\mathfrak{S}}\) implies

\[
\deg (Q _ {\mathfrak {S}}) \leqslant \deg (R _ {\mathfrak {S}}) \leqslant \deg (R)
\]

for all \(\mathfrak{S} \in \mathcal{B}\) ; whence immediately the required result by definition of \(\mathbf{Q}\) .
% semantic-fragments: legacy-input-seam
\relax{}
